\section{附录2：建模过程的创新点说明}

\subsection*{创新点一：三模型对比以实证手段验证线性假设}

\textbf{问题背景：}逻辑回归因其系数可解释性，在需要向业务方解释建模结论
时具有天然优势。但若特征间存在显著的非线性交互效应，线性假设将导致模型
欠拟合，预测精度明显低于非线性模型。

\textbf{创新做法：}本文并未简单地"因为需要解释性所以选用逻辑回归"，
而是通过逻辑回归、随机森林、XGBoost 三者的 AUC 对比（差值 $< 0.01$）
和 ROC 曲线叠合，实证地论证了该数据集上线性假设的充分性。
这一做法将模型选型的依据从主观偏好提升为可检验的实证结论。

\subsection*{创新点二：四层递进的探索性数据分析策略}

\textbf{问题背景：}传统的 EDA 通常仅凭某一类统计量（如均值差异或相关系数）
做出结论，缺乏多角度交叉验证。

\textbf{创新做法：}本文将 EDA 设计为"整体流失率 $\to$ 数值特征分布
$\to$ 分类特征流失率 $\to$ 相关性热力图"的四层递进结构：
第一层建立问题全貌认知，第二、三层从数值和分类两个角度独立探索，
第四层以相关系数的数值形式对前两层的定性发现进行量化锚定。
此外，问题一的 EDA 结论在问题二的 SHAP 分析和问题三的分层画像中
均有回扣（如在网时长的首位重要性与月付合同的高流失风险），
使全文形成"探索 $\to$ 验证 $\to$ 应用"的闭环。

\subsection*{创新点三：类别不平衡的全流程应对}

\textbf{问题背景：}客户流失数据中未流失客户占比 73.46\%，
流失客户仅占 26.54\%，接近 1:2.8 的比例。

\textbf{创新做法：}本文从探索到建模对类别不平衡进行了系统性应对：
\begin{enumerate}[nosep]
    \item 数据探索阶段：报告整体流失率，明确类别不平衡程度；
    \item 评估阶段：弃用对不平衡敏感的准确率，以 AUC 和 F1-score 为主要指标；
    \item 训练阶段：对逻辑回归设置 \texttt{class\_weight='balanced'}。
\end{enumerate}
这一贯穿全流程的应对，保证了评估指标真实反映模型预警能力，
而非给出"高准确率、零召回率"的虚假情况。

\subsection*{创新点四：分位数分层替代等距分层}

\textbf{问题背景：}将客户分为高/中/低风险是制定差异化策略的前提。
常见的等距切分（如 $[0,0.33)/[0.33,0.66)/[0.66,1]$）
在预测概率右偏分布下会导致中高风险层样本过于稀疏。

\textbf{创新做法：}采用 P50/P80 分位数作为切分点，
既保证了三个层级都有充足的样本量（50\%/30\%/20\%），
又使层级间的特征差异足够显著：高风险层的月付合同比例达 99.9\%、
光纤比例 88.9\%，与低风险层（19.5\% 和 22.9\%）形成鲜明对比。
